\section{Limitations and Future Work}
\label{appx:limitations_future}

\subsection{Goal Misgeneralization}
\label{misgeneralization}

One of the most common mistakes that \agentname{} makes during evaluation is to overgeneralize. For instance, if we prompt \agentname{} with a video of someone punching a cow, it may simply run to the nearest animal and punch that animal instead. This is related to the concept of goal misgeneralization \citep{misgeneralization}. Generalization can be helpful when the task we assign the agent is impossible to achieve from the current state and the agent instead performs a closely related action, but harmful when the task is achievable. We note two things: first, we believe the powerful generalization ability of \agentname{} probably comes from the MineCLIP embeddings and it especially improves the ability of \agentname{} to follow visual instructions when the exact items or blocks nearby are not available in the current environment, which is an extremely common scenario. Second, we notice that the tendency of the agent to misgeneralize decreases with scale. For example, with a model trained on less data, we find that asking the agent to look up and punch a tree to get a wooden log often resulted in the agent looking in the air and punching nothing; training the model on more data results in the agent first walking over to a nearby tree and looking up to get a wooden log. Future work should look to measure more closely the relationship between misgeneralization and scale.

\subsection{Random Selection of Hindsight Goals}

In this work, we choose to \textbf{randomly} select future episode timesteps as hindsight goals, rather than use a more sophisticated strategy. This is primarily due to the simplicity of the approach, but also to ensure a diverse and unbiased coverage of potential goals achievable within the short horizon. Future works can investigate the effects of alternative approaches that filter for semantically interesting timesteps as goals.

\subsection{Difference in Performance Between Text and Visual Goals}
In our experiments, we observed that \agentname{} often performed better when conditioned with visual goals compared to text goals converted through the prior. There are several potential factors that could contribute to this performance gap between text and visual conditioning. First, the prior model may not be accurately capturing the full meaning of the text prompt. Training the prior on more data or using a more powerful model architecture could potentially improve the quality of the sampled latent goals. Second, the visual goals can provide more precise demonstrations of the desired behavior. Text goals are inherently more ambiguous. Providing additional information such as the observation context to the prior and further prompt-engineering to make text prompts more detailed and less ambiguous, could help close this gap. Because text conditioning provides more flexibility and potential for generalization, closing the performance gap between text and visual conditioning is an important direction for future work.


\subsection{Challenges in Evaluating \agentname{}}
\label{appx:challenges_in_eval}

See \autoref{tab:other_tasks_table} for a non-exhaustive  list of tasks that the \agentname{} agent is able to achieve. It is worth mentioning that evaluating and describing the capabilities of open-ended generalist agents is an open research problem itself since capability depends strongly on preconditions, prompt engineering, and our own ability to come up with varied and challenging tasks. For example, there are many recent works on the evaluation of LLMs (e.g., \citep{helm, alpaca_eval, zheng2023judging}) which highlight these challenges.

That being said, there are a number of tasks which \agentname{} is unable to accomplish. As previously mentioned, long-horizon tasks such as obtaining a diamond or building a house are currently beyond the capability of \agentname{}. Further, \agentname{} also struggles with more complex crafting tasks like crafting an enchanting table, bookcase, or boat. Again, the virtually limitless and open-ended nature of tasks in Minecraft makes it very difficult to test generalist agents in this domain. We hope future works develop more sophisticated methods to evaluate the performance of generalist agents on short and long-horizon tasks (potentially through an extension of our MineCLIP evaluation method).

It is also worth noting that it is currently not possible to test \agentname{} or VPT \citep{baker2022vpt} in the MineDojo \citep{fan2022minedojo} environment, which is meant for generalist agent evaluation, since the action spaces are not compatible. We believe that bridging this gap could be greatly beneficial to the generalist agent community and we hope future works investigate this further.

\subsection{Towards Improved Long-Horizon Performance}
\agentname{} is a significant advancement in creating generative models of text-to-behavior, but it has several limitations. First, \agentname{} is mostly proficient at achieving short-horizon tasks while struggling on longer-horizon tasks like obtaining a diamond. Solving long-horizon tasks while taking actions using low-level mouse/keyboard controls is a very challenging and exciting research direction and, while prompt chaining is a promising approach for improving performance on complex tasks, more can be done in future work to improve performance.

One potential bottleneck is the fact that during packed hindsight relabeling, the hindsight goals are limited to at most 200 timesteps in the future ($\sim$10 seconds). Thus, tasks which require more than 200 timesteps to complete are technically out-of-distribution for \agentname{}. Although sampling hindsight goals from farther into the future could theoretically enhance long-horizon performance, our experiments in \Cref{appx:chunk_size} indicate that the performance tends to decrease if we increase this hyperparameter too much. We suspect that while increasing this hyperparameter may be able to improve long-horizon performance, it also increases noise and comes at the cost of reducing performance on short-horizon goals. Investigating whether it is possible to achieve a better tradeoff is an important avenue for future work. We also suspect that the long-horizon capabilities of \agentname{} could be improved through scaling or finetuning with reinforcement learning, or leveraging LLMs or VLMs to automatically provide prompt chains to the \agentname{} agent.


\subsection{Applying the \agentname{} Approach to Other Domains}

We designed \agentname{} for Minecraft due to the availability of two key ingredients: (1) a strong behavioral prior (VPT \citep{baker2022vpt}), and (2) a powerful visual-language model which maps text and video to a joint embedding space (MineCLIP \citep{fan2022minedojo}). However, the method used to create \agentname{} is not specific to the Minecraft domain. Given the rapid development of generative models, we expect that similar models to VPT and MineCLIP will become available in many other domains. As these models become available, future work could investigate the applicability of the \agentname{} approach to these other domains.